%% kapa: exam=1 round=24 year=2013 period=2 form=A subject=appraisal_law subject_ko=부동산관계법규 seq=31 number=71
%% source: Q-Net 감정평가사 기출문제 (cst00309) artlSeq=5073802 file=extracted/1008/5073802/2013년 제24회 감정평가사 1차시험 기출문제/2교시(A형).hwp
%% conversion: kordoc 4.24.51 -> tools/gen1.py (auto); review=proofread-sonnet; flags=none
\begin{question}{71}
건축법령상 건축위원회에 관한 설명으로 옳지 않은 것은?
\begin{choices}
  \item 국토교통부장관, 시ㆍ도지사 및 시장ㆍ군수ㆍ구청장은 각각 건축위원회를 두어야 한다.
  \item 시장ㆍ군수ㆍ구청장은 자신이 설치하는 건축위원회에 건축분쟁전문위원회를 둘 수 없다.
  \item 자치구의 경우에는 해당 자치구의 조례로 건축위원회의 조직ㆍ운영, 그 밖에 필요한 사항을 정한다.
  \item 건축분쟁전문위원회와 분야별 전문위원회의 심의등을 거친 사항은 건축위원회의 심의등을 거친 것으로 본다.
  \item 건축분쟁전문위원회와 분야별 전문위원회는 건축위원회가 정하는 사항에 대하여 심의등을 한다.
\end{choices}
\end{question}
\answer{3}
